% !TeX root = main.tex
\subsection{Overview of the Seven Test Cases}
The following seven scientific questions were used to evaluate the task organization and execution capabilities of Auto-ORCA Agent;
their intended capability coverage is summarized in Table~\ref{tab:agent-capability-design}.
Table~\ref{tab:task-overview} summarizes the methods actually used and the results obtained for each case.
T1--T5 and H1 were classified as successful calculations, whereas H2 remained incomplete, yielding a calculation success rate of $6/7=85.7\%$.
The numerical results and their scientific interpretation are presented below.

\begin{table}[H]
  \centering\small
  \caption{Methods actually used and results obtained for the seven test cases.
  Success indicates that an ORCA calculation terminated normally and produced the corresponding numerical results.
  }
  \label{tab:task-overview}
  \begin{tabularx}{\linewidth}{l>{\raggedright\arraybackslash}p{2.6cm}>{\raggedright\arraybackslash}p{3.1cm}>{\raggedright\arraybackslash}p{1.8cm}Y}
    \toprule
    ID & Question & Actual method & Status & Principal results\\
    \midrule
    T1 & Ethanol optimization and frequencies & r$^{2}$SCAN-3c & Success & Converged optimization; 21 positive frequencies and no imaginary frequencies\\
    T2 & Candidate energies for \textit{n}-hexane & r$^{2}$SCAN-3c single-point calculations & Success & Energy ranking of six candidate geometries\\
    T3 & Dihedral scan of \textit{n}-butane & r$^{2}$SCAN-3c constrained optimization & Success & Six converged points; actual angles remain unverified\\
    T4 & Vertical excitations of acetone & PBE0/def2-SVP, LR-TDDFT & Success & Five singlet excitation energies and oscillator strengths\\
    T5 & Environmental energy differences for ethanol & B3LYP/def2-TZVP and def2-SVP & Success & Gas-phase--CPCM water electronic energy differences with two basis sets\\
    H1 & Nuclear magnetic shielding of HCN & PBE0/def2-TZVP, GIAO & Success & Full shielding tensors for the C, N, and H nuclei\\
    H2 & Formaldehyde polarizability and rotational consistency & No actual calculation results & Incomplete & No polarizability values obtained\\
    \bottomrule
  \end{tabularx}
\end{table}

\subsection{T1: Geometry Optimization and Vibrational Frequencies of Ethanol}
\paragraph{Question}
Optimize the geometry of ethanol and use its vibrational frequencies to assess whether the resulting structure is a local minimum.

\paragraph{Method}
Geometry optimization and vibrational frequency calculations were performed using r$^{2}$SCAN-3c with its built-in def2-mTZVPP basis set for a neutral singlet in the gas phase.

\paragraph{Results}
\textbf{Successful calculation.} The geometry optimization converged, giving an optimized electronic energy of $-155.002345492244\,\Eh$.
The calculation yielded 21 positive vibrational frequencies, six zero-frequency translational and rotational modes after projection, and no imaginary frequencies.
All positive frequencies are listed in Table~\ref{tab:t1-frequencies}.
These results support assignment of the structure as a local minimum at the employed level of theory; they do not establish that it is the global minimum-energy conformer.
The optimized coordinates are retained in the XYZ file associated with the corresponding ORCA output.

\begin{table}[H]
  \centering\small
  \caption{The 21 positive vibrational frequencies of ethanol, in cm$^{-1}$.}
  \label{tab:t1-frequencies}
  \begin{tabular}{rrrrrrr}
    \toprule
    270.90 & 298.95 & 425.86 & 811.51 & 882.13 & 1052.13 & 1066.45\\
    1134.41 & 1282.25 & 1371.13 & 1396.18 & 1418.25 & 1483.33 & 1487.57\\
    1509.53 & 2969.59 & 3004.65 & 3061.83 & 3088.36 & 3105.66 & 3819.31\\
    \bottomrule
  \end{tabular}
\end{table}

\subsection{T2: Single-Point Electronic Energies of Candidate \textit{n}-Hexane Conformers}
\paragraph{Question}
Compare the electronic energies of candidate \textit{n}-hexane geometries and determine their relative energy ordering within the calculated set.

\paragraph{Method}
Single-point calculations were performed separately for six stored candidate geometries using r$^{2}$SCAN-3c for neutral singlets in the gas phase.
Relative energies were referenced to cand\_05.

\paragraph{Results}
\textbf{Successful calculation.} Electronic energies were obtained for all six candidates, as reported in Table~\ref{tab:t2-energies}.
Within this set, cand\_05 had the lowest energy, whereas cand\_03 was higher by 12.608656~kJ~mol$^{-1}$.
These are single-point results for candidate geometries; they do not demonstrate that all six conformers underwent DFT optimization,
nor do they establish the global minimum-energy conformer or a free-energy ordering.

\begin{table}[H]
  \centering\small
  \caption{Single-point electronic energies of six candidate \textit{n}-hexane geometries.}
  \label{tab:t2-energies}
  \begin{tabular}{lrr}
    \toprule
    Candidate & Electronic energy / $\Eh$ & Relative energy / kJ mol$^{-1}$\\
    \midrule
    cand\_05 & $-236.992066202851$ & 0.000000\\
    cand\_04 & $-236.991445556657$ & 1.629506\\
    cand\_01 & $-236.991444090308$ & 1.633356\\
    cand\_00 & $-236.991421448668$ & 1.692802\\
    cand\_06 & $-236.991245997872$ & 2.153448\\
    cand\_03 & $-236.987263819406$ & 12.608656\\
    \bottomrule
  \end{tabular}
\end{table}

\subsection{T3: Dihedral Scan of \textit{n}-Butane}
\paragraph{Question}
Obtain the electronic energies of \textit{n}-butane at different target dihedral angles through constrained optimization and examine the variation in torsional potential energy.

\paragraph{Method}
The main calculations employed r$^{2}$SCAN-3c for a neutral singlet in the gas phase, with constrained optimization at specified target dihedral angles.
An additional B3LYP/def2-TZVP calculation point is reported separately, together with the result of its geometry verification.

\paragraph{Results}
\textbf{Successful calculation.} Normal ORCA termination and optimization convergence were recorded for all six r$^{2}$SCAN-3c calculation points;
the results are given in Table~\ref{tab:t3-energies}.
The first column contains only the target-angle labels specified in the inputs.
The actual angles reported by the generated script conflict with these input labels and have not yet been confirmed through independent geometry verification.
Consequently, the energies are genuine computational outputs, but the labels cannot be treated as verified actual dihedral angles or used to establish a torsional barrier.
In particular, the high energy associated with the $+120^\circ$ label cannot be directly interpreted as a reliable physical barrier.

\begin{table}[H]
  \centering\small
  \caption{Six constrained optimization points for \textit{n}-butane.
  Angles denote input target labels; relative energies are the stored numerical values.
  }
  \label{tab:t3-energies}
  \begin{tabular}{rrr}
    \toprule
    Input target angle / $^\circ$ & Electronic energy / $\Eh$ & Relative energy / kJ mol$^{-1}$\\
    \midrule
    $-180$ & $-158.392470368263$ & 0.000000\\
    $-120$ & $-158.387146187943$ & 13.978634\\
    $-60$ & $-158.391526560891$ & 2.477966\\
    $0$ & $-158.383528626170$ & 23.476541\\
    $+60$ & $-158.391526702130$ & 2.477595\\
    $+120$ & $-158.358972460497$ & 87.948745\\
    \bottomrule
  \end{tabular}
\end{table}

The input target angle for the B3LYP/def2-TZVP calculation point was $-180^\circ$.
Independent geometry verification gave an actual angle of $179.999995786^\circ$, consistent with the target under angular periodicity;
the converged electronic energy was $-158.397118300397\,\Eh$.
Because this point was obtained using a different method, it is excluded from the relative energies in Table~\ref{tab:t3-energies}.

\subsection{T4: The First Five Singlet Vertical Excitations of Acetone}
\paragraph{Question}
Obtain the excitation energies, wavelengths, and oscillator strengths of the first five singlet vertical excitations of acetone.

\paragraph{Method}
Five singlet excited states were calculated at a fixed, prepared geometry using PBE0/def2-SVP for a neutral singlet in the gas phase,
with full linear-response TDDFT (LR-TDDFT; \code{TDA=false}).
Oscillator strengths were taken from the length-gauge output of ORCA.

\paragraph{Results}
\textbf{Successful calculation.} Convergence was recorded for both the SCF and response calculations; numerical results for the five excited states are listed in Table~\ref{tab:t4-excitations}.
The lowest singlet vertical excitation energy was 4.197061~eV, corresponding to a wavelength of 295.4~nm.

\begin{table}[H]
  \centering\small
  \caption{The first five singlet vertical excitations of acetone.
  Oscillator strengths $f$ are dimensionless.
  }
  \label{tab:t4-excitations}
  \begin{tabular}{crrr}
    \toprule
    Excited state & Excitation energy / eV & Wavelength / nm & Oscillator strength $f$\\
    \midrule
    $S_{1}$ & 4.197061 & 295.4 & 0.000007493\\
    $S_{2}$ & 7.464499 & 166.1 & 0.023781213\\
    $S_{3}$ & 8.406988 & 147.5 & 0.001899769\\
    $S_{4}$ & 8.572306 & 144.6 & 0.060276158\\
    $S_{5}$ & 8.640904 & 143.5 & 0.004121254\\
    \bottomrule
  \end{tabular}
\end{table}

\subsection{T5: Electronic Energy Differences between Gas-Phase Ethanol and Ethanol in CPCM Water}
\paragraph{Question}
Compare the electronic energies of ethanol at the same geometry in the gas phase and in a CPCM water environment.

\paragraph{Method}
Calculations employed B3LYP for a neutral singlet with the def2-TZVP and def2-SVP basis sets, respectively.
For each basis set, the gas-phase and aqueous-environment calculations used the same initial geometry; this geometry was not the optimized structure from T1.
The energy difference was defined as $\Delta E=E_{\mathrm{water}}-E_{\mathrm{gas}}$.

\paragraph{Results}
\textbf{Successful calculation.} Two ORCA calculations, one in the gas phase and one in CPCM water, were completed for each basis set.
Table~\ref{tab:t5-energies} reports the electronic energies, and Table~\ref{tab:t5-differences} reports their differences.
The negative signs indicate that, for the corresponding model and fixed geometry, the electronic energy in the aqueous environment was lower than that in the gas phase.
These quantities are \textbf{electronic energy differences, not hydration free energies}.

\begin{table}[H]
  \centering\small
  \caption{Electronic energies of ethanol in the gas phase and in CPCM water.}
  \label{tab:t5-energies}
  \begin{tabular}{lrr}
    \toprule
    Basis set & Gas-phase electronic energy / $\Eh$ & CPCM water electronic energy / $\Eh$\\
    \midrule
    def2-TZVP & $-155.002778990957$ & $-155.009740144896$\\
    def2-SVP & $-154.817054999208$ & $-154.823380641772$\\
    \bottomrule
  \end{tabular}
\end{table}

\begin{table}[H]
  \centering\small
  \caption{Electronic energy differences for ethanol in water relative to the gas phase.}
  \label{tab:t5-differences}
  \begin{tabular}{lrr}
    \toprule
    Basis set & $\Delta E$ / $\Eh$ & $\Delta E$ / kJ mol$^{-1}$\\
    \midrule
    def2-TZVP & $-0.006961153939$ & $-18.276507$\\
    def2-SVP & $-0.006325642564$ & $-16.607972$\\
    \bottomrule
  \end{tabular}
\end{table}

\subsection{H1: GIAO Nuclear Magnetic Shielding Tensors of HCN}
\paragraph{Question}
Obtain the nuclear magnetic shielding tensor and isotropic shielding of each atom in HCN.

\paragraph{Method}
Calculations were performed at a fixed initial geometry using PBE0/def2-TZVP and the GIAO method for a neutral singlet in the gas phase.
The atom order was C$_0$, N$_1$, and H$_2$.

\paragraph{Results}
\textbf{Successful calculation.} The total electronic energy was $-93.345070922944\,\Eh$.
The isotropic shielding values printed by ORCA are listed in Table~\ref{tab:h1-isotropic}; the full total shielding tensors are given in Eq.~\eqref{eq:h1-tensors}.
The tensor components are expressed in ppm, with both rows and columns ordered as X, Y, and Z; the three decimal places of the original output are retained.

\begin{table}[H]
  \centering\small
  \caption{Isotropic nuclear magnetic shielding values for HCN.}
  \label{tab:h1-isotropic}
  \begin{tabular}{cr}
    \toprule
    Atom & Isotropic shielding / ppm\\
    \midrule
    C$_0$ & 79.051\\
    N$_1$ & $-34.068$\\
    H$_2$ & 29.850\\
    \bottomrule
  \end{tabular}
\end{table}

\begin{equation}
  \begin{aligned}
    \bm{\sigma}_{\mathrm{C}_0}&=\begin{pmatrix}
      279.306 & -0.724 & -0.000 \\
      -1.638 & -21.077 & 0.000 \\
      -0.000 & 0.000 & -21.076
    \end{pmatrix},\\[0.8em]
    \bm{\sigma}_{\mathrm{N}_1}&=\begin{pmatrix}
      340.398 & -1.842 & -0.000 \\
      -0.044 & -221.291 & 0.000 \\
      0.000 & 0.000 & -221.309
    \end{pmatrix},\\[0.8em]
    \bm{\sigma}_{\mathrm{H}_2}&=\begin{pmatrix}
      39.712 & 0.055 & -0.000 \\
      -0.114 & 24.919 & -0.000 \\
      -0.000 & -0.000 & 24.920
    \end{pmatrix}.
  \end{aligned}
  \label{eq:h1-tensors}
\end{equation}
Isotropic shielding is one-third of the tensor trace; the tabulated values were taken directly from the summary printed by ORCA and were not recalculated from the rounded matrices.
These quantities have not been converted using a reference compound and therefore cannot be directly described as experimental chemical shifts.
H$_2$ denotes the index of the hydrogen atom and is unrelated to the task identifier H2 in the following subsection.

\subsection{H2: Static Polarizability and Rotational Consistency of Formaldehyde}
\paragraph{Question}
Obtain the static polarizability tensors and isotropic polarizabilities of formaldehyde in two orientations, and assess their physical consistency under rotation.

\paragraph{Method}
The geometries and rotation procedure were prepared, with the intention of calculating the polarizability in each orientation and comparing the results after coordinate transformation.

\paragraph{Results}
\textbf{Incomplete.} No actual ORCA polarizability values have been obtained.
Consequently, neither the polarizability tensors and isotropic polarizabilities in the two orientations nor a rotational consistency comparison is available.
The prepared geometries and rotation procedure do not constitute calculation results.

\paragraph{Analysis of the Failure}
The execution records indicate that interruption occurred primarily during task orchestration and data handoff before the calculation.
In the recorded attempts, an identity-description file was incorrectly bound as an XYZ geometry input, causing the rotation preprocessing script to terminate abnormally during coordinate parsing.
The identification and use of diagnostic information did not yield an effective repair; the corresponding runs terminated upon reaching either the model-call limit or the finite allowance for plan revisions.
